\subsection{ARCH-136-MOORE-PAIR: an exact one-disk high-degree repair}
\label{sec:arch136moorepair}

\paragraph{Statement.}
Fix a prime $p$, put $R=\mathbf Z_{(p)}$, and fix a positive weight
$m$.  Consider the strict differential PD algebra
\[
 D_p=\Lambda_R(a,c)\otimes_R\Gamma_R(b,t),
\]
with form degrees and weights
\[
 \begin{array}{c|cccc}
              &a&b&t&c\\\hline
 \text{form degree}&1&2&2p-2&2p-1\\
 \text{weight}&m&m&pm&pm
 \end{array}
\]
and differential
\[
 da=b,\quad db=0,\quad dc=b^{[p]},\quad
                    dt=pc-a b^{[p-1]}.                    \tag{136.MP1}
\]
Then $d^2=0$, the coefficient cohomology is $H^0(D_p)=R$, and
\[
 H^j(D_p)=0\quad(1\leq j\leq2p^2-1),
 \qquad H^{2p^2}(D_p)\simeq R/pR,                         \tag{136.MP2}
\]
where the last group has weight $p^2m$ and is generated by the
class of $b^{[p^2]}$.  Its full mod-$p$ cohomology algebra is
\[
 H^*(D_p\otimes_R\mathbf F_p)
  \simeq \Lambda_{\mathbf F_p}(\tau_1)\otimes
       \bigotimes_{j\geq2}
                   \mathbf F_p[z_j]/(z_j^p),             \tag{136.MP3}
\]
where $z_j$ is represented by $b^{[p^j]}$, has form degree $2p^j$
and weight $p^jm$, and $\tau_1$ is represented by
$c(b^{[p]})^{p-1}$, with form degree $2p^2-1$ and weight $p^2m$.
The initial disk $\Lambda(a)\otimes\Gamma(b)$ had its first
$p$-primary defect in degree $2p$; this pair of added generators
moves the first defect to degree $2p^2$ without changing the
algebra in form degrees below $2p-2$.

\paragraph{Proof.}
The special derivation rule gives
$d(a b^{[p-1]})=b b^{[p-1]}=p b^{[p]}$, so $d^2t=0$;
the other generator checks are immediate.  The free strictly
commutative graded PD algebra then carries this differential.
For example, $d^2$ on a divided power of $t$ is zero because
$d^2t=0$ and the odd element $dt$ squares to zero.  All formulas
preserve the stated weights.

Write $\overline D_p=D_p\otimes_R\mathbf F_p$.  The standard
factorial basis gives an algebra isomorphism
\[
 \Gamma_{\mathbf F_p}(b)
     \simeq\bigotimes_{j\geq0}\mathbf F_p[z_j]/(z_j^p),
                 \qquad z_j=b^{[p^j]}.                   \tag{136.MP4}
\]
Here is the integral coefficient check.  The coefficient expressing
$z_j^p$ as a multiple of $b^{[p^{j+1}]}$ has $p$-adic valuation
one, hence vanishes modulo $p$.  For an integer
$r=\sum_j e_jp^j$ with $0\leq e_j<p$, the coefficient of $b^{[r]}$
in $\prod_j z_j^{e_j}$ has valuation
\[
 v_p(r!)-\sum_j e_jv_p((p^j)!)=0.
\]
It is therefore a unit modulo $p$.  These monomials form the
factorial basis, proving (136.MP4), including injectivity.

Modulo $p$ the term $pc$ in $dt$ disappears.  Consequently
$\overline D_p$ is the tensor product of the three complexes
\[
 \begin{aligned}
 K_0&=\Lambda(a)\otimes\mathbf F_p[z_0]/(z_0^p)
                  \otimes\Gamma(t),
   &da&=z_0,\quad dt=-a b^{[p-1]},\\
 K_1&=\Lambda(c)\otimes\mathbf F_p[z_1]/(z_1^p),
   &dc&=z_1,\\
 K_{\geq2}&=\bigotimes_{j\geq2}
                          \mathbf F_p[z_j]/(z_j^p),
   &d&=0.
 \end{aligned}                                           \tag{136.MP5}
\]
All tensor products in this paragraph are over $\mathbf F_p$.
The coefficient relating $b^{[p-1]}$ to $z_0^{p-1}$ is the unit
$1/(p-1)!$.

The first complex contracts onto its constants.  Indeed in
$\Lambda(a)\otimes\mathbf F_p[z_0]/(z_0^p)$ the pairs
\[
                   a z_0^i\longmapsto z_0^{i+1}
                        \quad(0\leq i\leq p-2)
\]
are contractible.  The remaining vectors are $1$ and
$\tau_0=a z_0^{p-1}$.  Multiplication by $\tau_0$ annihilates
every nonconstant basis vector: a second $a$ squares to zero,
and a further positive power of $z_0$ reaches $z_0^p=0$.
It follows that tensoring each displayed contractible pair with
$\Gamma(t)$ leaves that pair as an actual contractible subcomplex;
the extra $dt$ terms are zero there.  On the remaining subcomplex,
\[
 d(t^{[r]})=-\frac1{(p-1)!}\tau_0 t^{[r-1]}
                                          \quad(r\geq1).
\]
These are unit two-term blocks, leaving only $\mathbf F_p$ in
degree zero.  This proves the asserted contraction of $K_0$ with
no spectral-sequence convergence issue.  It preserves weight.

The same elementary Koszul calculation, without a $t$ variable,
gives
\[
                 H^*(K_1)=\mathbf F_p\oplus\mathbf F_p\tau_1,
                 \qquad\tau_1=c z_1^{p-1}.
\]
Its square is zero.  Tensoring over the field with $K_{\geq2}$
proves (136.MP3).  Infinite tensor products cause no difficulty:
each degree uses only finitely many of the positive-degree
variables, and all elements have finite support.

Over $\mathbf Q$ replace $c$ by
\[
                         c'=c-\frac1p a b^{[p-1]}.
\]
Then $dc'=0$ and $dt=pc'$.  This is an invertible change of the
odd generator and identifies the rational algebra with the tensor
product of the odd unit disk $(a,b)$ and the even disk $(t,pc')$.
The former contracts over $\mathbf Q$ by dividing its factorial
block multipliers, and the latter contracts because $p$ is a unit.
Thus $D_p$ is rationally acyclic outside its constants.
Every form-degree term of $D_p$ is a finite free $R$-module, so
its positive cohomology is finitely generated $p$-primary torsion.

For any complex of free $R$-modules the coefficient exact sequence
gives
\[
 0\to H^j(D_p)/p\to H^j(\overline D_p)
                   \to H^{j+1}(D_p)[p]\to0.              \tag{136.MP6}
\]
By (136.MP3), the mod-$p$ groups vanish in degrees
$1,\ldots,2p^2-2$.  A nonzero finitely generated $p$-primary
module has both a nonzero quotient modulo $p$ and nonzero
$p$-torsion.  Therefore (136.MP6) implies the vanishing of all
integral groups in (136.MP2) below degree $2p^2$.

The closed element $b^{[p^2]}$ reduces to the nonzero class $z_2$
in (136.MP3).  Hence its integral class is nonzero and generates
$H^{2p^2}(D_p)$ modulo $p$; the middle group in (136.MP6) in that
degree has dimension one.  Finite generation and Nakayama's lemma
then show that this integral class generates the whole group.
Finally,
\[
 d(c b^{[p^2-p]})
       =\binom{p^2}{p}b^{[p^2]},\qquad
 \binom{p^2}{p}=p\binom{p^2-1}{p-1},                     \tag{136.MP7}
\]
and the last binomial factor is a $p$-adic unit.  The nonzero
cyclic group is therefore killed by $p$, proving that it is $R/pR$.
The class has its indicated weight.\hfill$\square$

\paragraph{Relation to the main construction.}
This proves exactly the one-disk phenomenon observed in the parent's
finite differential/mod-$p$ experiment, and goes beyond its tested
range.  It supplies a legitimate forward repair and quantifies its
next defect.  It does not yet define that repair naturally on a
simplicial free source or retain maps to several different theories.
The following calculation makes the repeated-cell issue precise.

\subsection{ARCH-136-MOORE-DUPLICATE: unrestricted duplicate repairs introduce lower torsion}

\paragraph{Statement and proof.}
Attach two copies $(c_1,t_1)$ and $(c_2,t_2)$ of (136.MP1) to the
same initial disk $(a,b)$.  Put
\[
                  c_\Delta=c_2-c_1,\qquad
                  t_\Delta=t_2-t_1.
\]
This invertible linear generator change is defined over $R$ and
preserves the free exterior and divided-power algebra structures.
The new differential is
\[
 dc_\Delta=0,\qquad dt_\Delta=p c_\Delta.
\]
Thus the double repair is the tensor product of the single repaired
algebra $D_p$ with
\[
 M_p=\Gamma_R(t_\Delta)\otimes\Lambda_R(c_\Delta),
       \quad |t_\Delta|=2p-2,\ |c_\Delta|=2p-1,
       \quad dt_\Delta=p c_\Delta.                       \tag{136.MD1}
\]
Its weight-$pm$ block in those degrees is $R\xrightarrow p R$,
so $[c_\Delta]\ne0$ in $H^{2p-1}(M_p)$ and has order $p$.
The augmentation of $D_p$ to $R$ is a chain algebra map and
splits the inclusion of $M_p$ into $D_p\otimes M_p$.
Consequently this nonzero class survives in the double repair.
The single pair was exact through $2p^2-1$, whereas the duplicate
already has a defect in degree $2p-1$.

Accordingly a construction attaching a free Moore pair for every
presentation of a torsion class cannot infer the one-pair range
from the individual calculation.  Identifying duplicates or
resolving their relations is additional mathematics, and preserving
maps to all targets must be checked at that step.  This counterexample
to that inference does not rule out a coherent repair functor or an
integral common source.\hfill$\square$

\subsection{ARCH-136-TAIL-MOORE: a single actual coefficient-tail repair lowers the first defect}

\paragraph{Statement.}
Let \(R=\mathbb Z_{(p)}\), \(S_R=\Gamma_R(T)\), and let
\(s\in S_R\) be homogeneous of positive weight \(e\), with
\(\bar s\neq0\) in \(\bar S=S_R/pS_R\). Start with the odd
unit disk \(da=b\), of form degrees one and two and common
weight \(m>0\). Attach just one free Moore pair for the actual
nonzero order-\(p\) class \(s\gamma_p(b)\):
\[
 \begin{gathered}
 A_s=S_R\otimes\Lambda_R(a,c)\otimes\Gamma_R(b,t),\\
 dc=s\gamma_p(b),\qquad
 dt=pc-s a\gamma_{p-1}(b),\\
 |c|=2p-1,\qquad |t|=2p-2,\qquad
 \operatorname{wt}(c)=\operatorname{wt}(t)=pm+e.
 \end{gathered}
 \tag{136.TM1}
\]
Then
\[
 H^j(A_s)=0\quad(1\leq j<2p-1),\qquad
 H^{2p-1}(A_s)
   \cong\operatorname{Ann}_{\bar S}(\bar s)
                       \langle pm+e\rangle.
 \tag{136.TM2}
\]
The last module is nonzero in arbitrarily high weights.
Thus a single repair of an actual nonzero coefficient-tail class,
with its minimal integer relation \(p\), already moves the first
\(p\)-defect from degree \(2p\) down to degree \(2p-1\).
No duplicate attachment or zero-cycle repair is used here.

\paragraph{Proof.}
The differential squares to zero because
\(d(a\gamma_{p-1}(b))=p\gamma_p(b)\).
The original disk has cohomology \(S_R/pS_R\) in form degree
\(2p\), generated by \(\gamma_p(b)\), and no positive
cohomology below that degree. Hence \(s\gamma_p(b)\) is indeed
nonzero of order \(p\).

The first new generator has degree \(2p-2\). Its coefficient
in a cycle of that degree must be zero, since its differential
has the independent \(c\)-coefficient \(p\) and \(S_R\) is
\(p\)-torsion-free. Thus the low vanishing in (136.TM2) follows
from that of the original disk.

In degree \(2p-1\), the only monomials over \(S_R\) are
\[
                  c,\qquad a\gamma_{p-1}(b),\qquad at.
\]
The coefficient of \(at\) in a cycle is zero: its differential
contains the independent monomial \(bt\), with coefficient one,
whereas the differentials of the other two displayed monomials
do not contain \(t\). Consequently every cycle is uniquely
\[
        f c-\frac{fs}{p}\,a\gamma_{p-1}(b),
                       \qquad f\in S_R,\quad fs\in pS_R.
 \tag{136.TM3}
\]
All boundaries in that degree come from multiples of \(t\);
the old degree-\((2p-2)\) monomial \(\gamma_{p-1}(b)\) is
closed. Under (136.TM3), the boundaries are exactly \(f\in pS_R\).
Reduction of \(f\) modulo \(p\) gives the claimed graded
\(S_R\)-module isomorphism.

Every positive-weight element of \(\bar S\) has \(p\)-th
power zero, by \(s^p=p!\gamma_p(s)\). Choose the least
\(r\geq2\) with \(\bar s^{r}=0\). Then
\(\bar s^{r-1}\neq0\) belongs to the annihilator.
It is homogeneous. A nonzero homogeneous element of \(\bar S\)
is a nonzero scalar times \(T^{[v]}\) for some \(v\).
For arbitrarily large powers \(p^a>v\), its product with
\(T^{[p^a]}\) is still nonzero: the coefficient
\(\binom{v+p^a}{v}\) is a \(p\)-unit by the no-carry
factorial calculation. Hence the annihilator is nonzero in
unbounded weights. This proves the last assertion.
\(\square\)

\paragraph{An explicit prime-power coefficient and the remaining choice.}
For \(s=T^{[N]}\), \(N=p^a\), set
\(u=p^{-1}\binom{pN}{N}\in R^\times\). The surviving class is
represented explicitly by
\[
 T^{[(p-1)N]}c-uT^{[pN]}a\gamma_{p-1}(b),
 \qquad
 p\bigl(T^{[(p-1)N]}c-uT^{[pN]}a\gamma_{p-1}(b)\bigr)
       =d\bigl(T^{[(p-1)N]}t\bigr).
 \tag{136.TM4}
\]
For the native repair \(s=1\), the annihilator in (136.TM2)
is zero, in agreement with ARCH-136-MOORE-PAIR. The difference
is therefore the coefficient-module relation, not an avoidable
choice to repair a zero class. Choosing actual cohomology classes
individually does not settle their \(S_R\)-relations. A coherent
nonfree or nonlinear repair could still behave differently;
no nonexistence statement for such a repair follows from this
one-pair calculation.
